\documentclass[main.tex]{subfiles}
\begin{document}

\chapter{Asymptotic Theory of the MLE and the Delta Method}\label{ch:mle-asimptotik}\chaptermark{MLE Asymptotics and the Delta Method}

The finite-sample properties of the MLE\index{maximum likelihood estimator mle@maximum likelihood estimator (MLE)} vary from model to model: sometimes it is biased, sometimes it has no closed form at all. In large samples, however, under regularity conditions\index{regularity conditions@regularity conditions}, the MLE shows universal and very strong behaviour: it is consistent, asymptotically normal\index{asymptotic normality@asymptotic normality}, and its asymptotic variance attains the Cramér--Rao\index{cramer rao lower bound crlb@Cramér--Rao lower bound (CRLB)} bound. This chapter covers these three results, the information inequality\index{information inequality@information inequality} that makes them possible, and the delta method\index{delta method@delta method}, which gives the asymptotic distribution of transformations of an estimator. These results are the theoretical basis of almost all standard errors and Wald confidence intervals\index{wald confidence interval@Wald confidence interval} reported in applied statistics.

\section{Definitions}

\begin{definition}[Asymptotic Normality and Asymptotic Variance]
If $\sqrt n(T_n-\theta)\dto\Normal(0,v(\theta))$, then $T_n$ is called \emph{asymptotically normal\index{asymptotic normality@asymptotic normality}} (AN) and $v(\theta)$ its \emph{asymptotic variance}; we write $T_n\approx\Normal(\theta,v(\theta)/n)$. The asymptotic variance need not equal $\lim n\Var(T_n)$ (the latter may not exist).
\end{definition}

\begin{definition}[Asymptotic Efficiency and Asymptotic Relative Efficiency]
In a regular model, $T_n$ is called \emph{asymptotically efficient} if $v(\theta)=1/I_1(\theta)$. For two AN estimators $T_n^{(1)},T_n^{(2)}$,
\[
\mathrm{ARE}(T^{(1)},T^{(2)})=\frac{v_2(\theta)}{v_1(\theta)}
\]
is the \emph{asymptotic relative efficiency\index{asymptotic relative efficiency@asymptotic relative efficiency}}. Interpretation: $T^{(2)}$ needs roughly $n/\mathrm{ARE}$ observations to reach the precision that $T^{(1)}$ achieves with $n$.
\end{definition}

\begin{definition}[Observed Information]
The random quantity $J_n(\theta)=-\ell''(\theta;\X)=-\sum_{i=1}^n\partial^2_\theta\ln f(X_i;\theta)$ is called the \emph{observed Fisher information\index{observed information@observed information}\index{fisher information@Fisher information}}. $\E_\theta J_n(\theta)=I_n(\theta)$ and, by the WLLN\index{law of large numbers@law of large numbers!weak wlln@weak (WLLN)}, $J_n(\theta)/n\pto I_1(\theta)$.
\end{definition}

\section{Consistency}

\begin{lemma}[The Information Inequality]\label{lem:bilgi-esitsizligi}
If $\theta\neq\theta_0$ and $\P_\theta\neq\P_{\theta_0}$ (identifiability), then
\[
\E_{\theta_0}\big[\ln f(X;\theta)\big]<\E_{\theta_0}\big[\ln f(X;\theta_0)\big].
\]
\end{lemma}
\begin{proof}
Since $\ln$ is strictly concave, Jensen\index{jensen s inequality@Jensen's inequality}'s inequality (Theorem~\ref{thm:jensen}) gives
\[
\E_{\theta_0}\ln\frac{f(X;\theta)}{f(X;\theta_0)}<\ln\E_{\theta_0}\frac{f(X;\theta)}{f(X;\theta_0)}=\ln\int f(x;\theta)\,dx=\ln1=0 .
\]
The strict inequality comes from the ratio not being a.s.\ constant (identifiability). The left-hand side is $-\mathrm{KL}(f_{\theta_0}\|f_\theta)$.
\end{proof}

\begin{theorem}[Consistency of the MLE \cite{wald49}]\label{thm:mle-tutarli}\index{maximum likelihood estimator mle@maximum likelihood estimator (MLE)}\index{consistency@consistency}
Let $X_i\iid f(\cdot;\theta_0)$ with an identifiable model. Suppose moreover that (i) $\Theta$ is compact, (ii) $\theta\mapsto\ln f(x;\theta)$ is continuous for every $x$, and (iii) $\E_{\theta_0}\sup_{\theta\in\Theta}|\ln f(X;\theta)|<\infty$. Then $\hat\theta_n\asto\theta_0$.
\end{theorem}
\begin{proof}
Let $M_n(\theta)=\frac1n\ell(\theta)$ and $M(\theta)=\E_{\theta_0}\ln f(X;\theta)$; by (iii), $M$ is finite everywhere. By Lemma~\ref{lem:bilgi-esitsizligi}, $M(\theta)<M(\theta_0)$ for $\theta\ne\theta_0$.

Fix $\varepsilon>0$ and let $K=\{\theta\in\Theta:|\theta-\theta_0|\ge\varepsilon\}$; $K$ is compact. For $\theta\in K$ and $r>0$ let $B(\theta,r)=\{\theta'\in\Theta:|\theta'-\theta|<r\}$ and
\[
m_r(x;\theta)=\sup_{\theta'\in B(\theta,r)}\ln f(x;\theta').
\]
By (ii), $m_r(x;\theta)\to\ln f(x;\theta)$ as $r\downarrow0$, and by (iii) $|m_r|\le\sup_{\Theta}|\ln f(x;\cdot)|$ is integrable. By dominated convergence, $\E_{\theta_0}m_r(X;\theta)\to M(\theta)<M(\theta_0)$; hence there is an $r_\theta>0$ with $\E_{\theta_0}m_{r_\theta}(X;\theta)<M(\theta_0)$. Let $B_1,\dots,B_m$ (centres $\theta^{(1)},\dots,\theta^{(m)}$) be a finite subcover of the open cover $\{B(\theta,r_\theta):\theta\in K\}$.

By the strong law of large numbers\index{law of large numbers@law of large numbers!strong slln@strong (SLLN)}, for each of the finitely many $j$ and for $\theta_0$, almost surely
\[
\frac1n\sum_{i=1}^n m_{r_j}(X_i;\theta^{(j)})\to\E_{\theta_0}m_{r_j}(X;\theta^{(j)})<M(\theta_0),\qquad M_n(\theta_0)\to M(\theta_0).
\]
On this event, for all sufficiently large $n$,
\[
\sup_{\theta\in K}M_n(\theta)\le\max_{j\le m}\frac1n\sum_{i=1}^nm_{r_j}(X_i;\theta^{(j)})<M_n(\theta_0)\le M_n(\hat\theta_n),
\]
i.e.\ $\hat\theta_n\notin K$: $|\hat\theta_n-\theta_0|<\varepsilon$. Intersecting over $\varepsilon=1/k$, $k\in\N$, gives $\hat\theta_n\asto\theta_0$. (For relaxing compactness and for the multidimensional case see \cite[Theorems~5.7 and 5.14]{vdv} and \cite[Chapter~17]{ferguson}.)
\end{proof}

\section{Asymptotic Normality}

\begin{theorem}[Asymptotic Normality of the MLE \cite{cramer46}]\label{thm:mle-an}\index{maximum likelihood estimator mle@maximum likelihood estimator (MLE)}
In addition to conditions (R1)--(R4) of Chapter~\ref{ch:cramer-rao}, let $\theta_0$ be an interior point, let $\hat\theta_n$ be a consistent root of the likelihood equation, and suppose $|\partial^3_\theta\ln f(x;\theta)|\le M(x)$ in a neighbourhood of $\theta_0$ with $\E_{\theta_0}M(X)<\infty$. Then
\[
\sqrt{n}(\hat{\theta}_{n} - \theta_0) \dto \Normal\left(0, \frac{1}{I_1(\theta_0)}\right).
\]
\end{theorem}
\begin{proof}
Expand the equation $\ell'(\hat\theta_n)=0$ around $\theta_0$: for some $\theta_n^*$ between $\hat\theta_n$ and $\theta_0$,
\[
0=\ell'(\theta_0)+(\hat\theta_n-\theta_0)\ell''(\theta_0)+\tfrac12(\hat\theta_n-\theta_0)^2\ell'''(\theta_n^*).
\]
Rearranging,
\[
\sqrt n(\hat\theta_n-\theta_0)=\frac{\frac{1}{\sqrt n}\ell'(\theta_0)}{-\frac1n\ell''(\theta_0)-\frac1{2n}(\hat\theta_n-\theta_0)\ell'''(\theta_n^*)} .
\]
\begin{itemize}
  \item \emph{Numerator:} $\ell'(\theta_0)=\sum s(\theta_0;X_i)$ is a sum of i.i.d.\ terms with mean $0$ and variance $I_1(\theta_0)$ (Lemma~\ref{lem:skor}). By the CLT\index{central limit theorem clt@central limit theorem (CLT)}, $\frac{1}{\sqrt n}\ell'(\theta_0)\dto\Normal(0,I_1(\theta_0))$.
  \item \emph{Denominator, first term:} by the WLLN\index{law of large numbers@law of large numbers!weak wlln@weak (WLLN)}, $-\frac1n\ell''(\theta_0)\pto I_1(\theta_0)$ (information identity\index{information identity@information identity}).
  \item \emph{Denominator, second term:} $|\frac1n\ell'''(\theta^*_n)|\le\frac1n\sum M(X_i)\pto\E M(X)<\infty$ and $\hat\theta_n-\theta_0\pto0$; the product $\pto0$.
\end{itemize}
By Slutsky\index{slutsky s theorem@Slutsky's theorem}'s theorem (Theorem~\ref{thm:slutsky}), $\sqrt n(\hat\theta_n-\theta_0)\dto\Normal(0,I_1)/I_1=\Normal(0,1/I_1(\theta_0))$.
\end{proof}

\begin{corollary}[The Wald Confidence Interval]\label{cor:wald}
If $I_1$ is continuous, then $I_1(\hat\theta_n)\pto I_1(\theta_0)$, and by Slutsky $\sqrt{nI_1(\hat\theta_n)}(\hat\theta_n-\theta_0)\dto\Normal(0,1)$. Hence
\[
\hat\theta_n\pm z_{\alpha/2}\big/\sqrt{nI_1(\hat\theta_n)}\qquad\text{or}\qquad\hat\theta_n\pm z_{\alpha/2}\big/\sqrt{J_n(\hat\theta_n)}
\]
are asymptotic $1-\alpha$ confidence intervals. In the multiparameter case, $\sqrt n(\hat{\bm\theta}_n-\bm\theta_0)\dto\Normal_k(\mathbf 0,I_1(\bm\theta_0)^{-1})$.
\end{corollary}
\begin{proof}
Since $\hat\theta_n\pto\theta_0$ and $I_1$ is continuous, Theorem~\ref{thm:cmt-p} gives $I_1(\hat\theta_n)\pto I_1(\theta_0)>0$. Since $Z_n=\sqrt{nI_1(\hat\theta_n)}(\hat\theta_n-\theta_0)=\sqrt{I_1(\hat\theta_n)}\cdot\sqrt n(\hat\theta_n-\theta_0)$, Slutsky gives $Z_n\dto\sqrt{I_1}\,\Normal(0,1/I_1)=\Normal(0,1)$, and
\[
\P_{\theta_0}\Big(\theta_0\in\hat\theta_n\pm z_{\alpha/2}/\sqrt{nI_1(\hat\theta_n)}\Big)=\P(|Z_n|\le z_{\alpha/2})\to1-\alpha
\]
(because $\pm z_{\alpha/2}$ are continuity points of $\Phi$). For the observed information\index{observed information@observed information} write $\frac1nJ_n(\hat\theta_n)=-\frac1n\ell''(\hat\theta_n)$; as in the proof of Theorem~\ref{thm:mle-an}, $|\frac1n\ell''(\hat\theta_n)-\frac1n\ell''(\theta_0)|\le|\hat\theta_n-\theta_0|\frac1n\sum M(X_i)\pto0$ and $-\frac1n\ell''(\theta_0)\pto I_1(\theta_0)$; the same argument applies. The multiparameter statement is the vector form of the proof of Theorem~\ref{thm:mle-an} (multivariate CLT for the score vector, Theorem~\ref{thm:cok-degiskenli-clt}).
\end{proof}

\section{The Delta Method}

\begin{theorem}[The Delta Method]\label{thm:delta}
Let $\sqrt{n}(T_n - \theta) \dto \Normal(0, \sigma^2)$ and let $g$ be differentiable at $\theta$ with $g'(\theta)\neq0$. Then
\[
\sqrt{n}\big(g(T_n) - g(\theta)\big) \dto \Normal\big(0, [g'(\theta)]^2 \sigma^2\big).
\]
\end{theorem}
\begin{proof}
Since $\sqrt n(T_n-\theta)=O_P(1)$, $T_n\pto\theta$. By the definition of the derivative, $g(t)-g(\theta)=[g'(\theta)+r(t)](t-\theta)$, where $r(t)\to0$ as $t\to\theta$ and $r(\theta)=0$. Hence
\[
\sqrt n\big(g(T_n)-g(\theta)\big)=\big[g'(\theta)+r(T_n)\big]\cdot\sqrt n(T_n-\theta).
\]
By the continuous mapping theorem\index{continuous mapping theorem@continuous mapping theorem} $r(T_n)\pto0$; by Slutsky\index{slutsky s theorem@Slutsky's theorem} the limit is $g'(\theta)\cdot\Normal(0,\sigma^2)$.
\end{proof}

\begin{theorem}[Second-Order Delta Method]\label{thm:delta2}
Under the conditions of Theorem~\ref{thm:delta}, if $g'(\theta)=0$ and $g''(\theta)\neq0$, then
\[
n\big(g(T_n)-g(\theta)\big)\dto\frac{\sigma^2g''(\theta)}{2}\,\chi^2_1 .
\]
\end{theorem}
\begin{proof}
Second-order Taylor: $n(g(T_n)-g(\theta))=\frac{g''(\theta)}{2}\big[\sqrt n(T_n-\theta)\big]^2+o_P(1)$; and $[\sqrt n(T_n-\theta)/\sigma]^2\dto\chi^2_1$ (continuous mapping).
\end{proof}

\begin{theorem}[Multivariate Delta Method]\label{thm:delta-cok}
If $\sqrt n(\mathbf T_n-\bm\theta)\dto\Normal_k(\mathbf 0,\Sigma)$ and $g:\R^k\to\R$ is differentiable at $\bm\theta$, then
$\sqrt n(g(\mathbf T_n)-g(\bm\theta))\dto\Normal\big(0,\nabla g(\bm\theta)^\top\Sigma\,\nabla g(\bm\theta)\big)$.
\end{theorem}
\begin{proof}
By the definition of differentiability, $g(\mathbf t)-g(\bm\theta)=\nabla g(\bm\theta)^\top(\mathbf t-\bm\theta)+\|\mathbf t-\bm\theta\|\,r(\mathbf t)$, where $r(\bm\theta)=0$ and $r$ is continuous at $\bm\theta$. Hence
\[
\sqrt n\big(g(\mathbf T_n)-g(\bm\theta)\big)=\nabla g(\bm\theta)^\top\sqrt n(\mathbf T_n-\bm\theta)+\big\|\sqrt n(\mathbf T_n-\bm\theta)\big\|\,r(\mathbf T_n).
\]
By the continuous mapping theorem (Theorem~\ref{thm:cmt}, linear $g$), the first term converges to $\nabla g^\top\mathbf Z$ with $\mathbf Z\sim\Normal_k(\mathbf 0,\Sigma)$, which is $\Normal(0,\nabla g^\top\Sigma\nabla g)$. In the second term, $\|\sqrt n(\mathbf T_n-\bm\theta)\|\dto\|\mathbf Z\|$, and hence $\mathbf T_n-\bm\theta=n^{-1/2}\cdot\sqrt n(\mathbf T_n-\bm\theta)\pto\mathbf 0$ (Slutsky); by continuity of $r$, $r(\mathbf T_n)\pto0$, and by Slutsky the product $\pto0$. Apply Slutsky once more.
\end{proof}

\begin{definition}[Variance-Stabilising Transformation]
If $\sqrt n(T_n-\theta)\dto\Normal(0,\sigma^2(\theta))$, then for the transformation $g(\theta)=\int\frac{c}{\sigma(\theta)}d\theta$, which satisfies $g'(\theta)=c/\sigma(\theta)$, we have $\sqrt n(g(T_n)-g(\theta))\dto\Normal(0,c^2)$: the asymptotic variance no longer depends on $\theta$.
\end{definition}

\section{Worked Examples}

\begin{example}[The Odds]\label{ex:odds}
$\hat p=\bar X$, $\sqrt n(\hat p-p)\dto\Normal(0,p(1-p))$. For the odds $\tau=g(p)=\frac{p}{1-p}$, $g'(p)=\frac{1}{(1-p)^2}$. By the delta method\index{delta method@delta method},
\[
\sqrt n(\hat\tau-\tau)\dto\Normal\Big(0,\frac{p(1-p)}{(1-p)^4}\Big)=\Normal\Big(0,\frac{p}{(1-p)^3}\Big).
\]
For the log-odds $\eta=\ln\frac p{1-p}$, $g'(p)=\frac{1}{p(1-p)}$ and the asymptotic variance is $\frac{1}{p(1-p)}$; this is the source of standard errors in logistic regression\index{logistic regression@logistic regression}.
\end{example}

\begin{example}[Variance-Stabilising Transformations]\label{ex:varyans-dengeleme}
\begin{itemize}
  \item \textbf{Poisson:} $\sigma^2(\lambda)=\lambda$; with $g(\lambda)=2\sqrt\lambda$, $\sqrt n(2\sqrt{\bar X}-2\sqrt\lambda)\dto\Normal(0,1)$.
  \item \textbf{Binomial:} $\sigma^2(p)=p(1-p)$; with $g(p)=\arcsin\sqrt p$ the asymptotic variance is $1/4$.
  \item \textbf{Correlation coefficient (Fisher's $z$):} for the bivariate normal, $\sqrt n(r-\rho)\dto\Normal(0,(1-\rho^2)^2)$; with $z=\tanh^{-1}r=\frac12\ln\frac{1+r}{1-r}$ the asymptotic variance is $1$.
\end{itemize}
These transformations both improve the coverage of confidence intervals and speed up the normal approximation.
\end{example}

\begin{example}[Asymptotic Relative Efficiency: Mean or Median?]\label{ex:are}
In a location family with symmetric density $f$, the asymptotic variance of the sample median is $\frac{1}{4f(\theta)^2}$ and that of the mean is $\sigma^2$.
\begin{center}
\begin{tabular}{@{}lccc@{}}
\toprule
Distribution & $v(\bar X)$ & $v(\text{median})$ & $\mathrm{ARE}(\text{median},\bar X)$ \\
\midrule
$\Normal(\theta,1)$ & $1$ & $\pi/2$ & $2/\pi\approx0.637$ \\
Laplace$(\theta,1)$ & $2$ & $1$ & $2$ \\
Cauchy$(\theta,1)$ & $\infty$ & $\pi^2/4$ & $\infty$ \\
\bottomrule
\end{tabular}
\end{center}
For normal data the median needs about $57\%$ more observations to match the efficiency of the mean; for the Laplace the median is twice as efficient (and is the MLE\index{maximum likelihood estimator mle@maximum likelihood estimator (MLE)}). For the Cauchy, the efficiency of the median relative to the MLE is $8/\pi^2\approx0.81$.
\end{example}

\begin{example}[Second-Order Delta: $p(1-p)$ at $p=1/2$]
$g(p)=p(1-p)$, $g'(1/2)=0$, $g''=-2$. Since $\sigma^2=1/4$ at $p=1/2$,
$n\big(\bar X(1-\bar X)-\tfrac14\big)\dto-\tfrac14\chi^2_1$. That is, at this point the estimator $\bar X(1-\bar X)$ converges at rate $n$ rather than $\sqrt n$, and its limit distribution is not normal but a (negatively scaled) chi-square.
\end{example}

\section{Computational Applications}

\subsection{The Delta Method for the Odds}

\pycode{code/ch12/c12_1_delta_method_odds.py}
\begin{lstlisting}[style=output]
Delta theoretical variance: 0.001749 | Empirical: 0.001764
Log-odds delta variance: 0.009524 | Empirical: 0.009570
\end{lstlisting}

\subsection{Asymptotic Normality of the MLE and Coverage of the Wald Interval}
For the Gamma shape parameter $\alpha$ (rate known, $\beta=1$), $I_1(\alpha)=\psi'(\alpha)$ (trigamma). The MLE\index{maximum likelihood estimator mle@maximum likelihood estimator (MLE)} solves $\psi(\hat\alpha)=\overline{\ln X}$.

\pycode{code/ch12/c12_2_mle_asymptotic_normality_wald.py}
\begin{lstlisting}[style=output]
n=  10 | bias=+0.0471 | n*Var=1.6058 (1/I=1.5505) | skewness(z)=+0.486 | Wald 95% coverage=0.9476
n=  50 | bias=+0.0113 | n*Var=1.5671 (1/I=1.5505) | skewness(z)=+0.219 | Wald 95% coverage=0.9501
n= 400 | bias=+0.0012 | n*Var=1.5378 (1/I=1.5505) | skewness(z)=+0.075 | Wald 95% coverage=0.9522
\end{lstlisting}

For small $n$ the positive bias and right skewness\index{skewness@skewness} are clear, and $n\Var(\hat\alpha)$ is above its asymptotic value; as $n$ grows, the bias vanishes at rate $O(1/n)$ and the skewness at rate $O(1/\sqrt n)$. In this model the \emph{total} coverage of the Wald interval\index{wald confidence interval@Wald confidence interval} is close to the target even at $n=10$, but this is because the errors in the two tails offset each other: because of the skewness, the interval misses too often on one side and too rarely on the other (Exercise: compute the two tail errors separately).

\subsection{ARE, Variance Stabilisation and the Second-Order Delta Method}

\pycode{code/ch12/c12_3_are_variance_stabilisation.py}
\begin{lstlisting}[style=output]
Normal  | n*Var(mean)=0.999 | n*Var(median)=1.561 | ARE(median, mean)=0.640
Laplace | n*Var(mean)=2.003 | n*Var(median)=1.156 | ARE(median, mean)=1.733
t_3     | n*Var(mean)=2.980 | n*Var(median)=1.862 | ARE(median, mean)=1.601
Poisson: Var(2*sqrt(Xbar)) * n  (target 1, independent of lambda)
  lambda=  1 | n*Var(Xbar)=  1.004 | n*Var(2 sqrt(Xbar))=1.006
  lambda=  5 | n*Var(Xbar)=  5.022 | n*Var(2 sqrt(Xbar))=1.005
  lambda= 50 | n*Var(Xbar)= 50.119 | n*Var(2 sqrt(Xbar))=1.002
Second-order delta: E[W]=-0.2487 (theory -0.25), Var[W]=0.1237 (theory 2/16=0.1250)
\end{lstlisting}

The normal and $t_3$ rows are very close to the theoretical ARE values ($0.637$ and $3\cdot4f(0)^2\approx1.62$): for the $t_3$ distribution (finite variance $=3$, but heavy-tailed) the median is more efficient than the mean. For the Laplace the empirical ARE ($\approx1.73$) is below the theoretical value $2$; since the Laplace density has a sharp (non-differentiable) peak, $n\Var(\text{median})$ approaches its asymptotic value $1$ only slowly, at rate $O(n^{-1/2})$. How well asymptotic results hold in finite samples should always be checked by simulation. The superiority of robust estimators for heavy-tailed data is a consequence of this table.

\section{Comments and Notes}

\begin{remark}[Hodges' Superefficient Estimator]
The estimator $T_n=\bar X\,\I\{|\bar X|\ge n^{-1/4}\}$ behaves like $\bar X$ in the $\Normal(\theta,1)$ model for $\theta\ne0$, while at $\theta=0$ its asymptotic variance is $0$: it seemingly ``beats'' the CRLB\index{cramer rao lower bound crlb@Cramér--Rao lower bound (CRLB)}. However, its risk deteriorates in a neighbourhood of $\theta=0$; Le Cam's theorem \cite{lecam53} shows that superefficiency\index{hodges superefficiency@Hodges' superefficiency} is possible only on a set of measure zero. This is why asymptotic efficiency\index{asymptotic relative efficiency@asymptotic relative efficiency} is defined together with local uniformity (regularity) (the Hájek--Le Cam convolution\index{convolution@convolution} theorem).
\end{remark}

\begin{remark}[Expected or Observed Information?]
Either $I_n(\hat\theta)$ or $J_n(\hat\theta)$ may be used in the Wald interval\index{wald confidence interval@Wald confidence interval}; both are asymptotically valid. The observed information\index{observed information@observed information} is usually more accurate from the point of view of conditional inference, because it reflects ancillary\index{ancillary statistic@ancillary statistic} information in the data, and it comes for free as a by-product of the Newton--Raphson\index{newton raphson method@Newton--Raphson method} algorithm.
\end{remark}

\begin{remark}[Limits of the Delta Method]
The delta method\index{delta method@delta method} is a local linearisation; if $g$ changes rapidly or the variance of $T_n$ is large, the approximation deteriorates. For example, $\hat\tau=\hat p/(1-\hat p)$ has infinite mean in finite samples because $\hat p=1$ has positive probability; yet its asymptotic distribution is well defined. In such cases the bootstrap (Chapter~\ref{ch:bootstrap}) or profile likelihood\index{profile likelihood@profile likelihood} intervals are preferred.
\end{remark}

% --- CHAPTER SUMMARY ---
\section{Chapter Summary}
\begin{summary}
\begin{itemize}
  \item By the information inequality (Lemma~\ref{lem:bilgi-esitsizligi}), the expected log-likelihood has its unique maximum at the true parameter. Consistency of the MLE (Theorem~\ref{thm:mle-tutarli}) follows from this.
  \item In regular models, $\sqrt n(\hat\theta_n-\theta_0)\dto\Normal(0,1/I_1(\theta_0))$ (Theorem~\ref{thm:mle-an}): the MLE attains the CRLB asymptotically.
  \item The Wald interval is built from the expected or the observed information (Corollary~\ref{cor:wald}). In small samples its coverage is sensitive to the parametrisation.
  \item The delta method (Theorem~\ref{thm:delta}) carries asymptotic normality over to differentiable transformations. When $g'(\theta)=0$ the second-order version (Theorem~\ref{thm:delta2}) is used, and in several dimensions the gradient version (Theorem~\ref{thm:delta-cok}). A variance-stabilising transformation makes the asymptotic variance constant.
  \item Hodges' superefficient estimator shows that asymptotic efficiency must be defined together with a local uniformity condition.
\end{itemize}
\end{summary}
% --- END OF CHAPTER SUMMARY ---

\paragraph{Further reading.} \cite[Chapter~10]{casella}, \cite[Chapters~3, 5 and 8]{vdv}, \cite[Chapters~17--18]{ferguson}.

\section{Exercises}

\begin{exercise}\label{exr:12-1}
Let $X_i\iid\text{Exp}(\lambda)$. Find the asymptotic distribution of $\hat\lambda=1/\bar X$ both by Theorem~\ref{thm:mle-an} and by the delta method. Show that the asymptotic variance of $\ln\hat\lambda$ does not depend on $\lambda$.
\hfill\hyperref[coz:12-1]{\footnotesize$\rightarrow$\,Solution}
\end{exercise}

\begin{exercise}\label{exr:12-2}
Let $(X_i,Y_i)$ be i.i.d.\ with $\E X=\mu_X\ne0$. Find the asymptotic variance of $\bar Y/\bar X$, the estimator of the ratio $\mu_Y/\mu_X$, by the multivariate delta method.
\hfill\hyperref[coz:12-2]{\footnotesize$\rightarrow$\,Solution}
\end{exercise}

\begin{exercise}\label{exr:12-3}
In the Cauchy location model, show (by simulation) that the one-step Newton--Raphson estimator starting from the median, $\tilde\theta=\hat m+\ell'(\hat m)/(n I_1)$ ($I_1=1/2$), is asymptotically efficient. Look up the general theory of such ``one-step'' estimators.
\hfill\hyperref[coz:12-3]{\footnotesize$\rightarrow$\,Solution}
\end{exercise}

\begin{exercise}\label{exr:12-4}
Using Fisher's $z$ transformation, construct an approximate $95\%$ confidence interval for $\rho$ from bivariate normal data. For $\rho=0.8$, $n=20$, compare its coverage by simulation with that of the direct interval $r\pm1.96(1-r^2)/\sqrt n$.
\hfill\hyperref[coz:12-4]{\footnotesize$\rightarrow$\,Solution}
\end{exercise}

\end{document}
